\documentclass[11pt,a4paper]{article}
\usepackage[margin=1in]{geometry}
\usepackage{amsmath,amssymb,amsthm}
\usepackage{algorithm}
\usepackage{algpseudocode}
\usepackage{booktabs}
\usepackage{hyperref}

\theoremstyle{definition}
\newtheorem{definition}{Definition}[section]
\newtheorem{theorem}{Theorem}[section]
\newtheorem{lemma}[theorem]{Lemma}
\newtheorem{remark}{Remark}[section]

\title{\textbf{Rigorous Mathematical Specification, Infinite-Width Feature Learning Limits, and Maximal Update Parametrization ($\mu\text{P}$)}}
\author{Team Scaibu \\ \small Email: \texttt{jaydeep.9664920749@gmail.com} \\ \small Architecture and Core Engine Team}
\date{October 2026}

\begin{document}
\maketitle

\begin{abstract}
This document presents the complete theoretical foundation and algorithmic specification of Maximal Update Parametrization ($\mu\text{P}$). In standard neural network parametrizations (Standard Parametrization / SP or Neural Tangent Kernel / NTK), the optimal learning rate shifts unpredictably as network width $n \to \infty$, necessitating expensive hyperparameter grid searches on large-scale models. $\mu\text{P}$ establishes the unique dimensional scaling of initialization variance ($\sigma^2 \propto 1/n$) and layer-wise learning rates ($\eta \propto 1/n$ for hidden and output projections, $\mathcal{O}(1)$ for input embeddings) that guarantees non-trivial feature learning and zero-shot hyperparameter transfer from small proxy models to massive foundation models. We derive the $\mu\text{P}$ scaling equations, formulate the complete algorithmic pseudo-code matching the reference implementation, prove the feature learning limit, step through a numerical calculation, and address implementation considerations.
\end{abstract}

\section{Notation and Mathematical Preliminaries}

Consider a deep network with base hidden width $n_0 \in \mathbb{N}_{\ge 1}$ scaled to target width $n \in \mathbb{N}_{\ge 1}$.

\begin{itemize}
    \item $n_0$: Reference base hidden dimension of small proxy model (e.g., $n_0 = 256$).
    \item $n$: Scaled target hidden dimension of large target model (e.g., $n = 4096$).
    \item $r \triangleq \frac{n}{n_0} \in \mathbb{R}_{> 0}$: Width expansion ratio.
    \item $\eta_0 > 0$: Optimal base learning rate tuned on the proxy model.
    \item $\eta(n)$: Scaled learning rate assigned to target model layers.
    \item $s_{\text{init}}(n)$: Multiplier applied to initial weight standard deviation $\sigma_0$.
    \item $\text{layer\_type} \in \{\text{input\_embedding}, \text{hidden\_weight}, \text{output\_head}\}$: Functional role in the computational graph.
\end{itemize}

\begin{table}[h]
\centering
\caption{Summary of Mathematical Symbols for $\mu\text{P}$}
\begin{tabular}{lll}
\toprule
\textbf{Symbol} & \textbf{Type / Domain} & \textbf{Description} \\
\midrule
$n_0$ & $\mathbb{N}_{\ge 1}$ & Proxy model width \\
$n$ & $\mathbb{N}_{\ge 1}$ & Target model width \\
$r$ & $\mathbb{R}_{> 0}$ & Width ratio $n / n_0$ \\
$\eta_0$ & $\mathbb{R}_{> 0}$ & Proxy tuned learning rate \\
$\eta$ & $\mathbb{R}_{> 0}$ & Scaled learning rate \\
$s_{\text{init}}$ & $\mathbb{R}_{> 0}$ & Initialization standard deviation multiplier \\
\bottomrule
\end{tabular}
\end{table}

\section{Problem Definition and Core Concepts}

\subsection{Intuition and Motivation: Zero-Shot Hyperparameter Transfer}
In Standard Parametrization (SP), tuning hyperparameters (learning rate, weight decay, warmup) on a small 100M parameter model does not transfer to a 70B parameter model: the optimal learning rate collapses or shifts by orders of magnitude. The Neural Tangent Kernel (NTK) limit maintains infinite-width stability but operates in the lazy-training regime where representations do not learn features. Yang et al. (2021) showed that $\mu\text{P}$ is the \textit{unique} scaling parameterization where:
\begin{enumerate}
    \item Every hidden layer activation update $\Delta h = h_t - h_0$ has coordinate magnitude $\Theta(1)$ as width $n \to \infty$.
    \item The optimal learning rate $\eta^*$ is invariant to width, allowing direct zero-shot transfer from small proxy models to massive target models without retraining.
\end{enumerate}

\subsection{Formal Mathematical Definition}
\begin{definition}[$\mu\text{P}$ Scaling Rules]
Given width expansion ratio $r = n / n_0$, the layer-wise learning rates and initialization multipliers are scaled as:
\begin{itemize}
    \item \textbf{Input Embeddings / Projections}:
    \begin{equation}
    \eta_{\text{input}} = \eta_0, \quad s_{\text{init}} = 1.0 \label{eq:mup_input}
    \end{equation}
    \item \textbf{Hidden Weights (Attention, MLP, Convolutions)}:
    \begin{equation}
    \eta_{\text{hidden}} = \frac{\eta_0}{r}, \quad s_{\text{init}} = \frac{1}{\sqrt{r}} \label{eq:mup_hidden}
    \end{equation}
    \item \textbf{Output Unembedding / Classifier Head}:
    \begin{equation}
    \eta_{\text{output}} = \frac{\eta_0}{r}, \quad s_{\text{init}} = \frac{1}{r} \label{eq:mup_output}
    \end{equation}
\end{itemize}
\end{definition}

\section{Algorithmic Specification}

The operational execution implemented in \texttt{NnAlgoMaximalUpdateParam} is presented below.

\begin{algorithm}[H]
\caption{Maximal Update Parametrization ($\mu\text{P}$) Scaling Computation}
\begin{algorithmic}[1]
\Require Base proxy width $n_0 \ge 1$, target width $n \ge 1$, base learning rate $\eta_0 > 0$
\Require Layer type $\in \{\text{input\_embedding}, \text{hidden\_weight}, \text{output\_head}\}$
\Ensure Scaled learning rate $\eta$, init standard deviation multiplier $s_{\text{init}}$, width ratio $r$
\If{$n_0 < 1$ \textbf{or} $n < 1$ \textbf{or} $\eta_0 \le 0$}
    \State \textbf{throw} PreconditionFailureException(``Invalid widths or base learning rate'')
\EndIf
\State $r \gets \text{float}(n) / \text{float}(n_0)$
\If{$\text{layer\_type} = \text{hidden\_weight}$}
    \State $\eta \gets \eta_0 / r$
    \State $s_{\text{init}} \gets 1.0 / \sqrt{r}$
\ElsIf{$\text{layer\_type} = \text{input\_embedding}$}
    \State $\eta \gets \eta_0$
    \State $s_{\text{init}} \gets 1.0$
\ElsIf{$\text{layer\_type} = \text{output\_head}$}
    \State $\eta \gets \eta_0 / r$
    \State $s_{\text{init}} \gets 1.0 / r$
\Else
    \State \textbf{throw} PreconditionFailureException(``Unrecognized layer type'')
\EndIf
\State \Return $\{\text{scaled\_lr}: \eta, \text{init\_std\_multiplier}: s_{\text{init}}, \text{width\_ratio}: r\}$
\end{algorithmic}
\end{algorithm}

\section{Theoretical Guarantees and Maximal Feature Learning}

\begin{theorem}[$\mu\text{P}$ Dynamical Stability (Yang et al., 2021)]
Under $\mu\text{P}$, for any layer $l$, the pre-activation change $\Delta z_l = z_l(t) - z_l(0)$ after one gradient step satisfies:
\begin{equation}
\Theta(1) \le \|\Delta z_l\|_\infty \le \mathcal{O}(1) \quad \text{as } n \to \infty
\end{equation}
guaranteeing maximum possible representation feature movement without activation blowup or collapse.
\end{theorem}

\begin{proof}
Let $h_l \in \mathbb{R}^n$ with $\|h_l\|_2 \sim \sqrt{n}$. The gradient with respect to hidden weights is $\nabla_W \mathcal{L} = \delta_{l+1} h_l^T$, where coordinate $\delta_{l+1, i} = \Theta(1/n)$ under output head scaling $1/r$. The update is $\Delta W = -\eta_{\text{hidden}} \nabla_W \mathcal{L} = -\frac{\eta_0}{n} (\delta_{l+1} h_l^T)$.
Computing the change in activation: $\Delta z_{l+1} = \Delta W h_l = -\frac{\eta_0}{n} \delta_{l+1} (h_l^T h_l)$. Since $h_l^T h_l = \|h_l\|_2^2 = \Theta(n)$, we obtain:
\begin{equation}
\Delta z_{l+1} = -\frac{\eta_0}{n} \delta_{l+1} \cdot \Theta(n) = -\eta_0 \delta_{l+1} \cdot \Theta(1) = \Theta(1)
\end{equation}
which is invariant to width $n$.
\end{proof}

\subsection{Limitations}
\begin{itemize}
    \item \textbf{Attention Softmax Scaling}: In multi-head self-attention under $\mu\text{P}$, the attention logit division factor must be scaled as $1 / d_k$ rather than $1 / \sqrt{d_k}$ to prevent logit explosion when head dimension scales.
\end{itemize}

\section{Complexity Analysis}

\subsection{Time and Space Complexity}
\begin{itemize}
    \item \textbf{Time}: Exactly $\mathcal{O}(1)$ scalar arithmetic operations.
    \item \textbf{Space}: $\mathcal{O}(1)$ memory.
\end{itemize}

\section{Worked Numerical Example}

Consider base proxy width $n_0 = 256$, target width $n = 4096$, tuned base learning rate $\eta_0 = 0.004$.

\begin{enumerate}
    \item \textbf{Width Expansion Ratio}:
    \begin{equation}
    r = \frac{4096}{256} = 16.0
    \end{equation}
    \item \textbf{Hidden Weight Projections}:
    \begin{align}
    \eta_{\text{hidden}} &= \frac{0.004}{16.0} = 0.000250 \\
    s_{\text{init}} &= \frac{1}{\sqrt{16.0}} = \frac{1}{4.0} = 0.250
    \end{align}
    \item \textbf{Output Head Projection}:
    \begin{align}
    \eta_{\text{output}} &= \frac{0.004}{16.0} = 0.000250 \\
    s_{\text{init}} &= \frac{1}{16.0} = 0.06250
    \end{align}
    \item \textbf{Input Embeddings}:
    \begin{equation}
    \eta_{\text{input}} = 0.0040, \quad s_{\text{init}} = 1.00
    \end{equation}
\end{enumerate}

\section{Numerical Considerations and Edge Cases}

\begin{itemize}
    \item \textbf{Adam vs. SGD Scaling in $\mu\text{P}$}: Under Adam, coordinate-wise division by $\sqrt{v_t}$ already cancels a factor of $1/\sqrt{n}$. The scaling formulas presented match Adam $\mu\text{P}$ standard formulations (e.g., in Microsoft \texttt{mup} library).
    \item \textbf{Multi-Head Dimension Independence}: If scaling width by adding more attention heads while keeping head dimension $d_{\text{head}}$ fixed, attention temperature remains standard $1/\sqrt{d_{\text{head}}}$.
\end{itemize}

\begin{thebibliography}{9}
\bibitem{yang2021} G.~Yang et al., ``Tensor Programs V: Tuning Large Neural Networks via Zero-Shot Hyperparameter Transfer,'' in \emph{Advances in Neural Information Processing Systems (NeurIPS)}, vol.~34, 2021.
\bibitem{yang2020} G.~Yang, ``Tensor Programs II: Neural Tangent Kernel for Any Architecture,'' \emph{arXiv preprint arXiv:2006.14548}, 2020.
\bibitem{dey2023} N.~Dey et al., ``Cerebras-GPT: Open compute-optimal language models trained on the Cerebras Wafer-Scale Cluster,'' \emph{arXiv preprint arXiv:2304.03208}, 2023.
\end{thebibliography}

\end{document}
